%!TeX root=MemoriaTFG.tex

\chapter{Introducció}
\label{cap:introduccio}

\section{La intel·ligència artificial en el joc}

Al llarg de la història, els jocs han servit com a banc de prova privilegiat per a la
intel·ligència artificial (\acsu{IA}). La raó és clara: un joc defineix un entorn tancat,
amb regles precises i un criteri d'èxit objectiu, la qual cosa permet mesurar el progrés
d'un agent de manera inequívoca. Els èxits en aquest àmbit han estat espectaculars.
L'any 1997, Deep Blue derrota el campió mundial d'escacs Garry Kasparov, un fita que
semblava inalcançable una dècada abans~\cite{campbell2002deepblue}. El 2016, AlphaGo supera
Lee Sedol al joc del Go, un problema considerat exponencialment més complex pels seus
\(10^{170}\) estats possibles~\cite{silver2016mastering}. El 2019, l'agent OpenAI Five
derrota equips professionals de Dota~2 en partides de cinc jugadors per banda, demostrant
la viabilitat de l'\acf{AR} en entorns multiagent de gran escala~\cite{berner2019dota}.

El motor comú de tots aquests avenços és el \emph{Deep Reinforcement Learning} (o
\acf{DRL}): la combinació de l'\ac{AR} clàssic amb xarxes neuronals profundes.
L'\ac{AR} permet a un agent aprendre comportaments òptims interactuant amb un entorn i
rebent recompenses; les xarxes profundes li permeten representar espais d'estats molt
complexos sense necessitat de definir les característiques rellevants manualment.
Junts, han fet possible que màquines sense cap coneixement inicial assoleixin nivells
sobrehumans en dominis on els humans tarden anys a dominar.

\section{El repte de la informació parcial}

Els èxits anteriors, però, es produeixen principalment en jocs d'\emph{informació
completa}: tots els jugadors veuen l'estat complet del joc en tot moment. Els escacs,
el Go i els jocs de taulell en general pertanyen a aquesta categoria. En canvi, molts dels
jocs que els humans troben més fascinants ---i estratègicament rics--- amaguen informació:
les cartes dels adversaris, les intencions futures, els recursos ocults.

Formalment, un joc d'informació parcial es modela com un \acf{POMDP}: l'agent
rep observacions parcials de l'estat real del món i ha d'actuar sense visibilitat
completa. Això introdueix un repte qualitatiu diferent, ja que l'estratègia òptima
ja no consisteix a calcular la millor jugada donada la posició, sinó a raonar sobre
la \emph{distribució de probabilitat} dels estats ocults, i a vegades a explotar
deliberadament la incertesa de l'adversari.

Els jocs de pòquer han estat el principal camp d'experimentació en informació parcial.
El 2017, Libratus guanya jugadors professionals de Texas Hold'em de dos jugadors
(\emph{heads-up})~\cite{brown2018libratus}. El 2019, Pluribus estén aquest resultat
al pòquer de sis jugadors, un entorn multiagent molt més complex on les tècniques
clàssiques de teoria de jocs no escalen directament~\cite{brown2019pluribus}.
Ambdós agents aproximen un \emph{equilibri de Nash}: una estratègia tal que cap jugador
pot millorar el seu resultat canviant unilateralment de política.

Malgrat aquests progressos, el camp és lluny d'estar tancat. Cada joc d'informació
parcial té la seva pròpia estructura d'observació, el seu espai d'accions i les seves
dinàmiques estratègiques. Un agent entrenat al pòquer no es pot transferir directament
a un altre joc de cartes: cal dissenyar l'entorn, definir l'espai d'estats i triar els
algorismes adequats a cada cas.

\section{El Truc: engany, parella i farol}

El Truc és un joc de cartes tradicional d'arrel catalana i valenciana, estès a
Llatinoamèrica i amb una comunitat activa a les Illes Balears i al País
Valencià~\cite{retruc}. Es juga per parelles (dos contra dos), amb l'objectiu de ser
el primer equip a arribar a 24 punts ---un \emph{cantó}--- en partides de millor de
tres o cinc cantons.

S'usa una baralla espanyola de 40 cartes de les quals s'eliminen els 2s, deixant
36 cartes amb quatre colls (Ors, Copes, Espases i Bastos). La jerarquia de les cartes
és no intuïtiva: la carta més alta no és l'as genèric, sinó l'As de Bastos (anomenat
\emph{L'Amo}), seguit del 10 d'Ors (\emph{Sa Madona}), l'As d'Espases i l'As de Bastos,
les set d'Espases i d'Ors, els tresos, la resta d'assos, les figures i finalment els
valors baixos.

Cada mà es divideix en dues batalles independents:

\begin{itemize}
  \item \textbf{L'Envit.} Abans de jugar cap carta, un jugador pot apostar sobre la
    força de la seva mà, calculada com la suma dels valors de dues cartes del mateix
    coll més 20 punts (les figures valen 0 en aquest còmput). L'escala d'apostes
    progressa des de \emph{Envit} (2 punts en joc) fins a \emph{La Falta} (tots els
    punts que li falten a l'equip guanyador per completar la cama). Si l'adversari
    no vol l'aposta, cedeix 1 punt a l'apostador.

  \item \textbf{El Truc.} En qualsevol moment, un jugador pot cantar \emph{Truc},
    posant 3 punts en joc. L'aposta pot escalar fins a \emph{Retruc} (6 pts),
    \emph{Quatre Val} (9 pts) i \emph{Joc Fora} (guanya tota la cama). L'objectiu
    és guanyar dues de les tres rondes (\emph{bassetges}) jugant la carta amb
    rang més alt en cada ronda.
\end{itemize}

La dimensió que fa del Truc un cas d'estudi especialment interessant no és, però,
el seu sistema de puntuació: és el \textbf{farol}. A diferència del pòquer, on el
farol és una acció aïllada, al Truc forma part de l'essència del joc. Un jugador
amb una mà mediocre pot cantar Truc amb convicció esperant que l'adversari es retiri,
cedint un punt sense veure les cartes. I alhora, les parelles poden comunicar
informació sobre les seves mans mitjançant un sistema de \emph{senyes} (gestos
facials codificats), sempre vigilant que els rivals no els deixin descoberts. Com
resumeix la saviesa popular del joc: \emph{``El Truc no és només matemàtiques i
memòria. És un joc de cares, de silencis i de mentides.''}

Des del punt de vista formal, el Truc és un \ac{POMDP} cooperatiu-competitiu:
l'agent ha de gestionar informació oculta (les cartes dels adversaris i de la parella),
coordinar-se amb un aliat sense comunicació verbal explícita, i decidir en cada moment
si enganyar o dir la veritat sobre la força de la seva mà. Aquesta combinació de
factors el converteix en un repte d'\ac{AR} substancialment diferent dels jocs
estudiats fins ara en la literatura.

\section{Objectius i estructura d'aquest treball}

L'objectiu d'aquest \acf{TFG} és desenvolupar i avaluar agents d'\acf{AR} capaços
de jugar al Truc de manera competitiva, partint de zero i sense cap coneixement
previ del joc. Concretament, el treball s'articula al voltant de tres objectius
principals:

\begin{enumerate}
  \item \textbf{Modelització de l'entorn.} Dissenyar un simulador del Truc que
    implementi totes les regles del joc (envit, truc, apostes, puntuació i senyes),
    i exposar-lo com un entorn estàndard compatible amb les biblioteques d'\ac{AR}
    existents (RLCard~\cite{zha2019rlcard} i Stable-Baselines3~\cite{raffin2021sb3}).

  \item \textbf{Comparació d'algorismes i tècniques d'entrenament.} Avaluar quatre
    algorismes ---\acf{DQN} i \acf{NFSP} via RLCard, i \ac{DQN} i \acf{PPO} via SB3---
    en condicions homogènies, i estudiar l'efecte del \emph{curriculum learning}
    (entrenar primer en mans individuals i després en partides senceres), de l'ús
    d'extractors de característiques convolucionals, de la memòria recurrent (GRU) i
    del \emph{self-play} sobre la convergència i el rendiment dels agents.

  \item \textbf{Avaluació.} Mesurar el rendiment de cada agent contra un agent de
    regles heurístiques i analitzar si els agents aprenen comportaments estratègics
    (gestió de l'envit, ús del farol al truc) o es limiten a polítiques reactives.
\end{enumerate}

Com a objectiu addicional, s'explora la variant per parelles del Truc, que afegeix
la dimensió cooperativa i la necessitat de coordinar-se amb un agent aliat.

El treball s'organitza de la manera següent. El capítol~\ref{cap:marc} presenta el
marc teòric: els fonaments de l'\ac{AR}, els algorismes emprats i l'estat de l'art
en jocs d'informació parcial. El capítol~\ref{cap:joc} descriu la lògica del joc i
l'arquitectura del simulador. El capítol~\ref{cap:entorns} detalla el disseny dels
entorns d'\ac{AR}. El capítol~\ref{cap:metodologia} explica la metodologia
d'entrenament per fases. El capítol~\ref{cap:resultats} presenta i discuteix els
resultats experimentals. Finalment, el capítol~\ref{cap:conclusions} recull les
conclusions i les línies de treball futur.
